\honest{Why Our Kind Can't Cooperate}{Eliezer Yudkowsky}{2009}

When I was young and made to attend, our synagogue held an annual appeal. The rabbi and the treasurer spoke about expenses, and the members called out their pledges from their seats.

Years later I ran an annual appeal of my own, online, for a nonprofit ``that may not be named,'' to an audience of atheists, libertarians, technophiles, science-fiction fans and programmers. \nb{The nonprofit was the Singularity Institute, now the Machine Intelligence Research Institute; in LessWrong's first months, AI and the Singularity were off-topic.} At once, people posted to the mailing lists their reasons for not giving. Some questioned the nonprofit's mission. Others had ideas for other sources of money; they did not offer to approach them. Yet donations came in right away, some with warm notes; one gift of \$111.11 came with the words ``One more hundred, one more ten, one more single, one more dime, and one more penny.'' But not one donor said so on the lists. Each donor, so far as they knew, was alone, and saw only arguments that they should not have given. It was as though, after the treasurer spoke, everyone not pledging stood to proclaim their reasons, while those pledging whispered. You might ask whether the objections were right; I ask you to hold that thought.

Someone I know asks why the Raelian flying-saucer cult can get tens of thousands of members, ``probably around 40,000,'' when we cannot get a thousand. \nb{Wikipedia says the movement ``claims tens of thousands of members''; the interview the post links could not be reached. Many Raelians call their creed an ``atheistic religion,'' so the contrast is not between believers and unbelievers.} Copying the cult by adding nonsense would be wrong, and might fail anyway, since there ``may be a hundred failed flying-saucer cults for every one that becomes famous,'' and the Dark Side needs skills you may lack; if you discuss your planned lies on the public Internet, you fail. Resenting a group you think inferior for having more followers also leads to the Dark Side. But if ``rationalists'' cannot coordinate as well as a flying-saucer cult, we should ask why.

Cults are held together by forces described in my earlier posts. The leader speaks, doubters keep quiet, and each member, seeing apparent unanimity, grows more confident: pluralistic ignorance. The unpersuaded leave, and those who remain agree more: evaporative cooling. Praise of the ideas feeds on itself, especially when criticism counts as treason, through the halo effect: an affective death spiral. Since we would not dirty our hands with these, it would seem to follow that the ``Light Side'' is always divided and weak, that atheists, libertarians, scientists and the rest will never act with the fanatic unity of radical Islam, and that ``the future inevitably belongs to the Dark.'' I state this only to reject it. Some ``Clash of Civilizations'' writers accept it with a sigh; ``I suppose'' they are signalling their sophistication. \nb{No writer is named.} I have always thought that a true rationalist ought to be effective in the real world, so surely unreason has its disadvantages. If current rationalist groups cannot support a project as well as one synagogue supports itself, finish that syllogism yourself.

My diagnosis is that it is dangerous to be half a rationalist. Teach someone only the fallacies they can find in others, and you make them stupider; this repeats ``Knowing About Biases Can Hurt People.'' Teach a group only individual rationality, and it cannot act together. I will write later about how rationalists might coordinate better.

The first of two forces against us is the ``culture of objections.'' At a conference, the questioners after a talk object to a logarithmic scale on slide 14, dispute a claim on slide 3, and offer an alternative hypothesis. Normal. Now picture questioners who say only that they agree with everything and ``You're awesome.'' You would flee in terror, as if Cthulhu had erupted from the podium. A group that tolerates disagreement but not agreement is also irrational; it hears only some honest thoughts. We are as uncomfortable together as cultists are apart, and reversed stupidity is not intelligence. An army in which every private knows strategy should not lose to one that follows orders; ``You definitely should not do worse'' with more knowledge; an uncoordinated mob gets slaughtered. A half-rationalist can easily do worse: in one experiment, politically opinionated students who knew more about the issues reacted less to contrary evidence, having more ammunition to argue against it. We are stuck in a valley of partial rationality, less coordinated than fundamentalists. A training program would teach both disagreeing and agreeing, comfort with dissent and with conformity: one day everyone dresses differently, another day in uniform. Practicing agreement and applause sounds like an evil cult, but why is only practicing disagreement acceptable? Are you never going to have to agree with the majority? \nb{For a single decision maker this is a standard result: extra information cannot lower the value of a decision, since it can be ignored.} Our culture rewards only heroic dissent; we signal intelligence and nonconformist membership by inventing clever objections, and ``Perhaps that is why'' our crowd stays marginalized. The title says ``Can't.'' \nb{In the year before the post, supporters of Ron Paul raised over \$4.2 million online in one day, and a British atheist bus campaign, after a failed pledge drive, raised over \pounds100,000 in four days. Such groups have coordinated giving at times, and have also failed to.}

The second force is shame of strong feeling. We still picture rationality as Spock's dispassion, or as cynicism, signalling sophistication by caring less than others. Would it not make you uncomfortable if a speaker said he cared so much about fighting aging that he would die for it? Yet decision theory contains no rule against caring. We should aim to feel the emotions that fit the facts, not to feel none. I quote ``That which can be destroyed by the truth should be'' and add my own ``That which the truth nourishes should thrive.'' \nb{The first line is P.~C. Hodgell's, as Yudkowsky's ``Twelve Virtues of Rationality'' says. The pairing restates ``Feeling Rational'' (2007).} Some things are worth dying for. The taboo on emotional language in science papers may help the facts fight it out, but it need not apply everywhere; poetic appeals for things that need doing deserve applause. We need the strength to resist ungrounded appeals and the strength to be moved by grounded ones.

The synagogue had it right. No one went row by row asking Mr.~Schwartz how much he would give. People announced their pledges simply, without drama, and that ``probably'' helped others overcome weakness of will. \nb{Experiments support this: in a field experiment on charitable giving, Shang and Croson (2009) found that telling new donors what someone else had given raised their gifts.} Those with objections voiced them at another time. At the least, support and objections should both be heard. A wall of objections is the ``mirror image'' of what ``Dark Side groups'' do.
